%------------------------------------------------------------------------------%
\documentclass[fleqn,utf8,aspectratio=169,12pt]{beamer}

\usepackage{calculo_varias_variaveis-1}

\title{Derivadas Parciais de Ordem Superior}

%------------------------------------------------------------------------------%
\begin{document}

\addtitle

%------------------------------------------------------------------------------%
\section{Derivadas Parciais de Segunda Ordem}

%------------------------------------------------------------------------------%
\begin{frame}{Conceito}
  Derivada de uma derivada de uma função
\end{frame}

%------------------------------------------------------------------------------%
\begin{frame}{Notações}
\begin{columns}
\column{0.5\linewidth}
  Derivadas de segunda ordem
  \[
    \frac{\partial^2 f}{\partial x^2} \qquad\qquad\qquad
    f_{xx}
  \]

  \vspace{\baselineskip}
  \[
    \frac{\partial^2 f}{\partial y^2} \qquad\qquad\qquad
    f_{yy}
  \]
\column{0.5\linewidth}
  \pause
  Derivadas de segunda ordem mistas
  \[
    \frac{\partial^2 f}{\partial x \partial y} \qquad\qquad\qquad
    f_{yx}
  \]

  \vspace{\baselineskip}
  \[
    \frac{\partial^2 f}{\partial y \partial x} \qquad\qquad\qquad
    f_{xy}
  \]
\end{columns}
\end{frame}

%------------------------------------------------------------------------------%
\begin{frame}{Atenção Para a Ordem das Derivadas Mistas}
  \[
    \frac{\partial^2 f}{{\color{magenta} \partial x} {\color{blue} \partial y}}
    = \frac{\partial}{{\color{magenta} \partial x}}
      \left(\,\frac{\partial f}{{\color{blue} \partial y}}\,\right)
  \]

  \vspace{\baselineskip}
  \[
    f_{{\color{blue} x}{\color{magenta} y}} = \big(\,f_{\color{blue} x}\,\big)_{\color{magenta} y}
  \]
\end{frame}

%------------------------------------------------------------------------------%
\include{F-03-Derivadas_parciais_ordem_superior-ex1}
\include{F-03-Derivadas_parciais_ordem_superior-ex2}

%------------------------------------------------------------------------------%
\section{Teorema das Derivadas Mistas}

%------------------------------------------------------------------------------%
\begin{frame}{Teorema das Derivadas Mistas}
  Se a função \(f(x, y)\) e suas derivadas parciais
  \, $f_\emph{x}$,\, $f_\emph{y}$,\, $f_{\emph{xy}}$\, e \,$f_{\emph{yx}}$

  \pause
  forem definidas em uma vizinhança do ponto \((a, b)\)

  \pause
  e \emph{todas forem contínuas}, então

  \pause
  \[
    \DM{f}{x}{y} = \DM{f}{y}{x}
  \]
\end{frame}

%------------------------------------------------------------------------------%
\include{F-03-Derivadas_parciais_ordem_superior-ex3}

%------------------------------------------------------------------------------%
\section{Derivadas Parciais de Ordem Superior}

%------------------------------------------------------------------------------%
\begin{frame}{Notações}
  \begin{columns}
    \column{0.5\linewidth}
    Derivadas terceiras
    \[
      \frac{\partial^3 f}{\partial x^3} \qquad\qquad\qquad
      f_{xxx}
    \]

    \vspace{\baselineskip}
    \[
      \frac{\partial^3 f}{\partial y\partial x^2} \qquad\qquad\qquad
      f_{xxy}
    \]
    \column{0.5\linewidth}
    Derivadas quartas
    \[
      \frac{\partial^4 f}{\partial x^4} \qquad\qquad\qquad
      f_{xxx}
    \]

    \vspace{\baselineskip}
    \[
      \frac{\partial^4 f}{\partial y^2 \partial x^2} \qquad\qquad\qquad
      f_{xxyy}
    \]
  \end{columns}
\end{frame}

%------------------------------------------------------------------------------%
\begin{frame}{Ordem de Derivação}
  Desde que todas as derivadas envolvidas sejam contínuas

  \pause
  a ordem de derivação é irrelevante
\end{frame}

%------------------------------------------------------------------------------%
\include{F-03-Derivadas_parciais_ordem_superior-ex4}

%------------------------------------------------------------------------------%
\section{Lista Mínima}

%------------------------------------------------------------------------------%
\begin{frame}{Lista Mínima}
  Cálculo Vol. 2 do Thomas 12\textsuperscript{a} ed. -- Seção 14.3
  \begin{enumerate}
    \item Estudar o texto da seção
    \item Resolver os exercícios: 43-45, 51-54, 57, 61
  \end{enumerate}

  \vfill
  Atenção: A prova é baseada no livro, não nas apresentações
\end{frame}

%------------------------------------------------------------------------------%
\end{document}
%------------------------------------------------------------------------------%
